The bug report is right that the earlier solution did **not** prove C3 as stated. The correct repair is not to pretend that the missing implications follow, but to state precisely what follows from A5 and what does not.

Below is an improved, logically justified solution.

---

## Improved solution

\[
\boxed{\text{C3, as stated, is not derivable from A5 alone.}}
\]

Let the four documented cases in A5 be denoted by

\[
\alpha=\text{AlphaFold2},\qquad
d=\text{AI drug discovery},\qquad
g=\text{GNoME},\qquad
\ell=\text{Lean4 AI proof assistance}.
\]

Clause (ii) of C3 asserts that the time-to-discovery speedup is at least \(2\times\) in **each** documented case, and exceeds \(100\times\) in at least one case. Thus, if \(S_x\) denotes the time-to-discovery speedup in case \(x\), clause (ii) requires

\[
S_\alpha\ge 2,\qquad S_d\ge 2,\qquad S_g\ge 2,\qquad S_\ell\ge 2,
\]

and also

\[
\max\{S_\alpha,S_d,S_g,S_\ell\}>100.
\]

A5 supplies enough information to prove this only for the drug-discovery case. In that case, the AI-assisted time is \(18\) months, while the historical human-only average is \(4\)--\(6\) years. Since

\[
4\text{ years}=48\text{ months},
\]

the smallest speedup supported by A5 is

\[
S_d\ge \frac{48}{18}=\frac{8}{3}>2.
\]

So A5 proves

\[
S_d>2.
\]

However, A5 does **not** provide corresponding time-to-discovery data for the other three cases:

1. For AlphaFold2, A5 says the protein-structure-prediction problem had been open for \(50\) years and was solved to near-experimental accuracy, but it gives no AI elapsed time \(T_\alpha^{AI}\). Therefore no numerical speedup

   \[
   S_\alpha=\frac{T_\alpha^{\text{human}}}{T_\alpha^{AI}}
   \]

   can be computed from A5 alone.

2. For GNoME, A5 gives a discovery-count ratio exceeding \(100\times\). It does **not** give a time-to-discovery speedup. The statement

   \[
   \text{GNoME discovered }2.2\text{ million stable crystal structures}
   \]

   compared with approximately \(20{,}000\) prior human experimental discoveries supports a scale ratio, not a time ratio. Moreover, because the prior value is stated as approximately \(20{,}000\), one should not treat

   \[
   \frac{2{,}200{,}000}{20{,}000}=110
   \]

   as an exact derivation. The rigorous conclusion is the one explicitly stated in A5:

   \[
   \text{GNoME ratio}>100\times.
   \]

3. For Lean4 AI proof assistance, A5 says that AI-assisted proof search produced novel proof strategies for open combinatorics problems and that such tasks previously required years of specialist effort. But A5 gives no AI-assisted elapsed time. Therefore no numerical speedup \(S_\ell\) can be computed.

Thus A5 proves one explicit time-speedup claim,

\[
S_d>2,
\]

and one explicit discovery-scale claim,

\[
R_g>100,
\]

where \(R_g\) is a count ratio, not a time-to-discovery speedup. It does not prove

\[
S_x\ge 2
\]

for every documented case \(x\), nor does it prove that some time-speedup \(S_x\) exceeds \(100\).

Therefore clause (ii) of C3 does not follow from A5.

Since a theorem requiring clauses (i), (ii), and (iii) cannot be proved if clause (ii) is not derivable, C3 as stated is not provable from A5 alone.

---

## About clause (i)

Clause (i) also requires care. It says that in at least three independent scientific domains, AI achieved results

\[
\text{strictly beyond pre-AI human capability, not merely faster}.
\]

A5 supports a weaker, record-based statement: in several cases, AI produced results that had not previously been achieved by the human scientific record described in A5. For example:

- AlphaFold2 solved a problem that A5 says had been open for \(50\) years.
- GNoME produced \(2.2\) million novel stable crystal structures and exceeded the prior human experimental corpus by more than \(100\times\).
- Lean4 AI produced novel proof strategies for open combinatorics problems.

But “not previously achieved by humans” is weaker than “strictly beyond pre-AI human capability.” The latter could mean that pre-AI humans could not have achieved the result at all, even with more time or resources. A5 does not establish that stronger modal claim.

So clause (i) is provable only under an additional interpretive convention, such as:

\[
\text{“beyond pre-AI human capability” means “beyond the pre-AI human record described in A5.”}
\]

Under that operational reading, AlphaFold2, GNoME, and Lean4 AI can support clause (i). Under the stronger literal reading, A5 is insufficient.

---

## A corrected theorem actually provable from A5

A theorem that is genuinely derivable from A5 is the following.

```tex
\begin{theorem}[Scientific Force Multiplication, Derivable Form]
  Under Assumption~A5, AI systems constitute a cross-domain scientific
  force multiplier in the following precise sense:
  \begin{enumerate}
    \item[(a)] A5 documents AI-driven scientific gains in four explicitly
      named problem areas: protein-structure prediction, drug discovery,
      stable crystal-structure discovery, and formal proof search in
      combinatorics.
    \item[(b)] At least three of the documented cases produced results not
      merely described as faster executions of an already completed prior
      human result: AlphaFold2 solved a problem open for $50$ years, GNoME
      produced novel stable crystal structures exceeding the prior human
      experimental corpus by a stated ratio of $>100\times$, and Lean4 AI
      produced novel proof strategies for open combinatorics problems.
    \item[(c)] A5 documents a time-to-discovery speedup greater than
      $2\times$ in the drug-discovery case.
    \item[(d)] A5 documents a discovery-count multiplication greater than
      $100\times$ in the GNoME case.
  \end{enumerate}
\end{theorem}
```

### Proof

By A5, the four documented AI-driven results concern the following explicitly named problem areas:

\[
\text{protein-structure prediction},
\]

\[
\text{drug discovery},
\]

\[
\text{stable crystal-structure discovery},
\]

and

\[
\text{formal proof search in combinatorics}.
\]

These are four distinct problem descriptions appearing directly in A5. Therefore the documented gains are not confined to a single problem area.

For AlphaFold2, A5 states that protein-structure prediction was a problem open for \(50\) years and that AlphaFold2 solved it to near-experimental accuracy. A result that solves a previously open problem is not merely a faster repetition of an already completed solution, because the prior human record described in A5 did not contain such a solution.

For GNoME, A5 states that it discovered \(2.2\) million novel stable crystal structures and that this surpassed the total of all prior human experimental discovery by a stated ratio of more than \(100\times\). Since the structures are described as novel, the result is not merely a faster reproduction of already known structures. Also, because A5 itself states the ratio as

\[
>100\times,
\]

we may conclude directly that the discovery-count multiplication exceeds \(100\times\). We need not treat the approximate value \(20{,}000\) as exact.

For Lean4 AI proof assistance, A5 states that AI-assisted proof search produced novel proof strategies for open combinatorics problems. Since the strategies are described as novel and the problems as open, this again is not merely a faster execution of an already recorded prior solution.

For the drug-discovery case, A5 states that an AI-designed drug candidate reached IND approval in \(18\) months, whereas the historical human-only average was \(4\)--\(6\) years. The smallest stated human-only baseline is \(4\) years, or \(48\) months. Hence the minimum speedup supported by A5 is

\[
\frac{48}{18}=\frac{8}{3}>2.
\]

Thus A5 proves a time-to-discovery speedup greater than \(2\times\) in the drug-discovery case.

Combining these points, A5 establishes cross-domain scientific force multiplication, including one rigorously quantified time-speedup exceeding \(2\times\) and one rigorously stated discovery-count multiplication exceeding \(100\times\). However, without additional time-to-discovery data, it does not establish C3(ii) as originally written. \(\square\)

---

So the corrected conclusion is:

\[
\boxed{\text{A5 supports a strong cross-domain force-multiplication claim, but not C3 exactly as stated.}}
\]